\chapter{CONCLUSIONES Y RECOMENDACIONES}
    \section{Conclusiones}
        La capacidad excepcional del modelo para armonizar las características espectrales y espaciales de las imágenes MSS con las TM se demostró de manera concluyente. La implementación del Transformer Visual y del Perceptrón Multicapa dentro de una arquitectura GAN aseguró una alineación precisa, como se evidencia en la alta fidelidad visual y la correspondencia espectral de las imágenes resultantes.

        La aplicación de técnicas avanzadas de aprendizaje profundo condujo a una mejora notable en la corrección geométrica de las imágenes MSS, lo que permitió una superposición casi perfecta con las imágenes TM. Este avance en la corrección geométrica es un paso significativo hacia la homogeneización y estandarización de datos satelitales para análisis a largo plazo.

        La investigación mostró que es posible complementar imágenes MSS con bandas espectrales no originalmente presentes a través de un modelo de aprendizaje profundo. Este desarrollo no solo enriquece las imágenes MSS sino que también potencia su utilidad para aplicaciones que requieren un mayor detalle espectral.

        La estratégica selección de una región diversa para la validación del modelo y la comparación con modelos previos confirmaron la robustez y la generalización del modelo desarrollado. Los resultados subrayaron la necesidad de adaptar el modelo a las variaciones topográficas y espectrales específicas de cada zona de interés.

    \section{Recomendaciones}
        Para futuras investigaciones en el campo de la armonización de imágenes satelitales, se recomienda enfocarse en el desarrollo y la experimentación con nuevas arquitecturas de redes neuronales, como los Transformers y GANs avanzados, para mejorar la calidad de la armonización espectral y la generación de bandas virtuales. Este enfoque debería ir acompañado de un esfuerzo por diversificar y ampliar los conjuntos de datos utilizados, incluyendo imágenes de una variedad más amplia de condiciones ambientales y geográficas, lo que permitirá evaluar la robustez y la generalizabilidad de los modelos propuestos.

        Además, se sugiere la creación de herramientas automatizadas y fácilmente integrables con sistemas de información geográfica existentes, lo que facilitaría la aplicación práctica de los modelos de armonización en diversos campos como la agricultura, el urbanismo y el estudio del cambio climático. La publicación de código fuente y conjuntos de datos en plataformas de acceso abierto potenciará la colaboración interdisciplinaria y acelerará la adopción de estas tecnologías avanzadas en la comunidad científica y profesional, contribuyendo así al avance en el monitoreo global y análisis detallado del cambio en la superficie terrestre.